\section{Background and Motivation}

\subsection{Hypervisor--OS Coupling as an Architecture Problem}

A hypervisor mediates between hardware and one or more guest virtual machines
(VMs), but it cannot do so in isolation: it requires a substrate to allocate
memory, manage execution contexts, route interrupts, access devices, and
coordinate initialization. How this substrate is provided varies widely across
designs, yet in nearly every case the hypervisor is architecturally fused to a
particular substrate choice. This fusion is not accidental---it reflects design
decisions about where to draw the boundary between virtualization logic and OS
services---but it has consequences that are familiar to software engineering:
high coupling, low portability, and duplicated effort whenever the
virtualization layer must be retargeted.

KVM is the canonical example of full integration. The Kernel-based Virtual
Machine is not a standalone program but a Linux subsystem that exposes
virtualization through a \code{/dev/kvm} character device and a set of
ioctls~\cite{kivity2007kvm}. Its vCPU scheduling, memory management, interrupt
routing, and device handling are inseparable from the Linux kernel, which gives
KVM access to Linux's mature drivers and scheduler but makes it unthinkable
outside Linux. Firecracker, the lightweight VMM behind AWS Lambda, is built atop
KVM and therefore transitively inherits the same substrate
dependency~\cite{agache2020firecracker}. At the other extreme,
static-partitioning hypervisors such as Bao~\cite{martins2020bao} and
Jailhouse~\cite{ramsauer2017jailhouse} run with a much smaller trusted computing
base, but they are no less bound to their substrate: Jailhouse relies on a
booted Linux instance to initialize hardware and then partitions it, while Bao
assumes a bare-metal boot and memory model. ACRN targets IoT consolidation and
supports a service-VM model that is itself tightly coupled to a specific guest
configuration~\cite{li2019acrn}. Even research efforts that rethink the
hypervisor--kernel boundary, such as DuVisor's user-level hypervisor via
delegated virtualization~\cite{chen2022duvisor} or the Out-of-Hypervisor
approach to nested virtualization~\cite{bitchebe2022ooh}, relocate the coupling
rather than eliminate it: each still presupposes a particular substrate
architecture.

The practical consequence is that virtualization support is not a portable
capability but a per-OS investment. When a new OS---a Rust-based kernel such as
Asterinas~\cite{peng2025asterinas, peng2024framekernel}, a unikernel-style base
such as ArceOS, or an experimental research kernel---wishes to host VMs, it must
either reimplement a hypervisor, port an existing one by invasive surgery on
both sides, or do without virtualization entirely. The literature on
safety-critical and embedded virtualization confirms this fragmentation:
comparative studies of Jailhouse, Xen Dom0-less, Bao, and seL4 CAmkES treat each
hypervisor as an independent system whose internal coupling to its substrate is
a given, not a variable~\cite{martins2020bao, lozano2023comprehensive}. None of
these systems asks whether the \emph{same} virtualization core could be reused
across substrates.

\subsection{Two Dimensions of Decoupling}

Portability in virtualization is usually discussed along two axes. The first is
\emph{guest portability}: a hypervisor should run different guest operating
systems. The second is \emph{hardware portability}: a hypervisor should span
instruction-set architectures and platform configurations, as evidenced by work
on Arm and RISC-V virtualization~\cite{sa2023cva6}. Both axes have received
extensive attention. A third axis, however, is largely missing:
\emph{substrate-OS portability}, i.e., whether the hypervisor itself can be
reused across different host operating systems that provide its substrate
services.

We further observe that substrate-OS portability has two independent sides, and
a complete decoupling story must address both. The \emph{substrate-facing side}
concerns how the hypervisor obtains OS services below it: memory, CPUs,
interrupts, timers, devices, and synchronization. The \emph{guest-facing side}
concerns how guests and their toolchains interact with the hypervisor above it.
On the guest side, the KVM ioctl ABI has become a de facto standard: a large
ecosystem of VMMs, device backends, and guest drivers assumes a KVM-compatible
interface. If a portable hypervisor exposes this ABI, then unmodified guests and
toolchains work regardless of which substrate hosts the hypervisor. Decoupling
only the substrate side while ignoring the guest side would still force guests
to be rewritten per hypervisor; decoupling only the guest side would still force
the hypervisor to be rewritten per OS. Our work targets both sides
simultaneously.

\subsection{A Software-Engineering Perspective}

From a software-architecture standpoint, the hypervisor--OS boundary is an
instance of a classic problem: separating a reusable component from the platform
services it consumes through an explicit, minimal contract. Abstraction and
modularity are foundational principles for managing complexity, enabling reuse,
and bounding the impact of change~\cite{shaw2025abstractions, wan2023software}.
In practice, however, ``portability'' is frequently an aspiration rather than a
verified property, and abstraction boundaries that look clean on paper often
leak platform-specific semantics under stress~\cite{mince2023illusion}. This
tension is acute for hypervisors because the services they need---physical
memory ownership, interrupt-context execution, address-translation control---are
precisely the services where OS designs diverge most.

This makes the hypervisor--OS interface a demanding contract to design. It is
not a hardware abstraction layer that hides device registers; it is a service
contract that must expose enough low-level control for virtualization while
hiding the accidental structure of each substrate. If the contract is too
narrow, the hypervisor cannot map guest memory, inject interrupts, or drive
devices. If it mirrors one OS's abstractions too closely, portability
evaporates. If it papers over semantic differences, the same hypervisor code may
appear portable while silently relying on assumptions that hold on only one
substrate. A useful contract must therefore be small, explicit, and faithful to
the semantics that virtualization actually requires. Our central claim is that
such a contract exists, is feasible to implement, and that the cost it imposes
is modest relative to re-implementing a hypervisor per OS.

\subsection{Why Now: Emerging Substrates}

The motivation for tackling substrate-OS portability now is concrete. A new
generation of OSes is being built in memory-safe languages and with
non-traditional architectures: Asterinas adopts a framekernel design with a
small, sound trusted computing base while preserving Linux ABI
compatibility~\cite{peng2025asterinas, peng2024framekernel}; library-OS and
unikernel approaches such as Unikraft~\cite{kuenzer2021unikraft} and
LibrettOS~\cite{nikolaev2020librettos} demonstrate that OS construction can be
decomposed into reusable primitives; and unikernel isolation continues to be
explored for edge and embedded targets~\cite{kaiser2024unikernels,
ueda2024mewz}. Each of these substrates could benefit from
virtualization---for compatibility, testing, isolation, or cloud-native
execution---but none should have to absorb the cost of building a hypervisor
from scratch. If the hypervisor can be treated as a portable component that
consumes a defined service contract, then emerging OSes gain virtualization by
implementing a thin adaptation layer rather than a full virtualization stack.
This is the gap our work fills.
